\subsection{非分裂情形}

命题 14. 设 \(k'\) 是 \(k\) 的扩张，\(\mathfrak{g}' = \mathfrak{g} \otimes_k k'\) 。设 \(\mathfrak{m}\) 是 \(\mathfrak{g}\) 的子代数，\(\mathfrak{m}' = \mathfrak{m} \otimes_k k'\) 。若 \(\mathfrak{m}'\) 是 \(\mathfrak{g}'\) 的抛物（相应地 Borel）子代数，则 \(\mathfrak{m}\) 是 \(\mathfrak{g}\) 的抛物（相应地 Borel）子代数。

由命题 8 与 11，只需证明 m 包含 g 的一个分裂嘉当子代数即可。设 b 是 g 的一个 Borel 子代数。则 \(b' = b \otimes_{k} k'\) 是 \(\mathfrak{g}'\) 的 Borel 子代数，故 \(\mathfrak{m}' \cap \mathfrak{b}'\) 包含 \(\mathfrak{g}'\) 的一个嘉当子代数（命题 10）。设 \(\mathfrak{t}\) 是 \(\mathfrak{m} \cap \mathfrak{b}\) 的一个嘉当子代数。则 \(\mathfrak{t} \otimes_k k'\) 是 \(\mathfrak{m}' \cap \mathfrak{b}'\) 的、从而是 \(\mathfrak{g}'\) 的嘉当子代数（第七章 §3 第 3 节 命题 3）。于是，\(\mathfrak{t}\) 是 \(\mathfrak{g}\) 的嘉当子代数，且由于它含于 \(\mathfrak{b}\) ，它是分裂的。

定义 3. 设 \(\mathfrak{a}\) 是半单（或更一般地，约化）李代数，\(\bar{k}\) 是 \(k\) 的一个代数闭包。子代数 \(\mathfrak{m}\) （\(\mathfrak{a}\) 的）称为抛物的（相应地 Borel 的），如果 \(\mathfrak{m} \otimes_k \bar{k}\) 是 \(\mathfrak{a} \otimes_k \bar{k}\) 的抛物（相应地 Borel）子代数。

若 \(\mathfrak{a}\) 是可分裂的，则命题 14 表明这个定义与定义 2（相应地与定义 1）等价。

命题 15. 设 \(\mathfrak{a}\) 是约化李代数，\(k'\) 是 \(k\) 的扩张，\(\mathfrak{m}\) 是 \(\mathfrak{a}\) 的子代数。则 \(\mathfrak{m}\) 是 \(\mathfrak{a}\) 的抛物（相应地 Borel）子代数当且仅当 \(\mathfrak{m} \otimes_k k'\) 是 \(\mathfrak{a} \otimes_k k'\) 的抛物（相应地 Borel）子代数。

这由命题 14 立即得出。

